\section{Transport des opérateurs de masse par CKM}
\label{sec:transporte-ckm-operadores-rev11}

L'incidence N42 fixe le sextuplet \((4,5,7,6,8,3)\) et la matrice rationnelle
\[
 M_{\rm CKM}=
 \begin{pmatrix}
 2&-5&28\\
 0&7&-25/2\\
 1&-20&-2/3
 \end{pmatrix},
 \qquad \det M_{\rm CKM}=-\frac{3857}{6}.
\]
Les coordonnées angulaires sont construites par
\[
 \begin{pmatrix}\theta_{12}\\\theta_{23}\\\theta_{13}\end{pmatrix}
 =M_{\rm CKM}
 \begin{pmatrix}A\\ C^*/6\\ (180/\pi)\Delta_4\end{pmatrix},
 \qquad \delta_{\rm CP}=9A.
\]
Cette transformation appartient à la clôture mathématique de HMT. Dans sa réalisation physique, soient \(B_u\) et \(B_d\) les opérateurs bidirectionnels déjà construits dans les fibres de quarks de types up et down. Le changement de base agit selon
\begin{equation}
 B_d^{(u)}=V_{\rm CKM}B_dV_{\rm CKM}^*,
 \qquad
 R_{\rm act}^{B_d^{(u)}}
 =V_{\rm CKM}R_{\rm act}^{B_d}V_{\rm CKM}^*.
 \label{eq:transporte-ckm-kappa-rev11}
\end{equation}
CKM transporte les valeurs propres entre bases et conserve le spectre et le déterminant. Le contrôle de la violation CP prend la forme de l'identité
\[
 \left|\det[B_u,V_{\rm CKM}B_dV_{\rm CKM}^*]\right|
 =2|J|\prod_{i<j}|\kappa_i^u-\kappa_j^u|
       \prod_{a<b}|\kappa_a^d-\kappa_b^d|.
\]
La matrice de mélange et le lecteur de route remplissent donc des fonctions complémentaires : le lecteur détermine l'opérateur spectral et CKM détermine sa représentation dans la base de saveur.

\section{Repères géométriques de rang \(3\) pour le secteur leptonique}
\label{sec:frames-pmns-rev11}

La construction leptonique utilise l'involution diagonale de signe
\[
 S_{\rm sign}=\operatorname{diag}(1,1,1,-1,-1,-1,1,1,1,-1,-1,-1),
 \qquad P_-=(I-S_{\rm sign})/2,
\]
et le projecteur icosaédrique canonique \(P_3\) de rang trois. La compression
\[
 C_-=P_-P_3P_-\big|_{\operatorname{im}P_-}
\]
possède le spectre
\[
 \left\{\frac{5+\sqrt5}{10},\frac12,
 \frac{5-\sqrt5}{10},0,0,0\right\}.
\]
Les trois valeurs propres positives sont simples. Leurs projecteurs spectraux, appliqués à l'indicatrice de l'hexade négative orientée et normalisés après projection par \(P_3\), construisent un repère orthonormal chargé \(F_\ell\) qui satisfait
\[
 F_\ell^*F_\ell=I_3,
 \qquad F_\ell F_\ell^*=P_3.
\]

Pour le secteur des neutrinos, soient \(H_1,H_2,H_3\) les trois hexades positives orientées et
\[
 X=(P_3\mathbf1_{H_1},P_3\mathbf1_{H_2},P_3\mathbf1_{H_3}).
\]
Sa matrice de Gram est définie positive et satisfait
\[
 \det(X^*X)=\frac{5-\sqrt5}{40}>0.
\]
Par conséquent,
\[
 F_\nu=X(X^*X)^{-1/2},
 \qquad F_\nu^*F_\nu=I_3,
 \qquad F_\nu F_\nu^*=P_3.
\]
Le transport réel
\begin{equation}
 O=F_\ell^*F_\nu\in SO(3)
 \label{eq:transporte-frames-pmns-rev11}
\end{equation}
s'obtient sans masses ni angles PMNS.

\HMTClaim{MD-R-PMNS-RANKONE-UNITARY}
\begin{theorem}[Unitarité de la phase dodécaphasique minimale]
\label{thm:unitaria-interna-pmns-rev11} Soit \(b_K\in\operatorname{im}P_3\) le vecteur normalisé du témoin de Witt, \(Q_K=|b_K\rangle\langle b_K|\), et
\[
 \Phi_L=I+(\zeta_{12}-1)Q_K,
 \qquad \zeta_{12}=e^{i\pi/6}.
\]
Alors
\begin{equation}
 U_{\rm HMT}=F_\ell^*\Phi_LF_\nu
 \label{eq:unitaria-interna-pmns-rev11}
\end{equation}
est une matrice unitaire indépendante de valeurs métrologiques cibles. L'inversion d'orientation conjugue \(U_{\rm HMT}\).
\end{theorem}

\begin{proof}
\(Q_K\) est un projecteur orthogonal, donc \(\Phi_L\) est unitaire. Les deux repères sont des isométries de même image \(\operatorname{im}P_3\). Par conséquent
\[
 U_{\rm HMT}^*U_{\rm HMT}
 =F_\nu^*\Phi_L^*P_3\Phi_LF_\nu=I_3.
\]
Le changement \(\pi/6\mapsto-\pi/6\) produit \(\Phi_L^*\), et toutes les autres primitives sont réelles.
\end{proof}

\subsection{Contrôle négatif de rang un et transport PMNS de rang complet}
\label{sec:no-go-pmns-rango-uno-rev11} \HMTClaim{MD-R-PMNS-RANKONE-PHYSICAL-NOGO}

Une fois la construction figée, la comparaison dans l'ordre déclaré produit
\[
 (\theta_{12},\theta_{23},\theta_{13})
 =(19.6807^\circ,56.4970^\circ,29.6224^\circ),
 \qquad J=-0.0426621.
\]
La meilleure des 36 permutations discrètes laisse \(\theta_{13}\simeq17.02^\circ\). La phase de rang un engendre une matrice unitaire interne, mais conserve un stabilisateur \(U(2)\) et est exclue comme réalisation physique de PMNS.

L'exclusion du porteur de rang un ne laisse pas la réalisation PMNS indéterminée. Soient \(M_\ell\) et \(M_\nu\) les opérateurs chargé et neutre produits par la Loi générale des masses, et soient \(F_\ell,F_\nu\) leurs repères spectraux orthonormaux complets sur l'image commune \(\operatorname{im}P_3\). Le transport leptonique de rang complet est alors construit par
\[
 U_{\rm PMNS}:=F_\ell^*F_\nu,
 \qquad
 U_{\rm PMNS}^*U_{\rm PMNS}=I_3.
\]
Pour des spectres simples, cette sortie est déterminée à des réajustements diagonaux de phase et à l'ordre spectral déclaré près ; en présence de dégénérescences, la sortie exacte est la classe double de repères correspondante. Le noyau du registre constitue un certificat fini supplémentaire de cette réalisation, et non son générateur ou une condition encore requise pour son existence. Pour les fibres chargées \(L_a\) et neutriniques \(N_b\), sa forme agrégée est
\[
 G_{ab}=\frac{1}{\sqrt{|L_a||N_b|}}
 \sum_{c\in L_a}\sum_{d\in N_b}K_{\rm reg}(c,d),
 \qquad U_{\rm int}=\operatorname{polar}(G),
\]
avec naturalité sous l'orientation. Lorsque les entrées explicites du registre vérifient \(\det G\ne0\), son facteur polaire fournit un certificat alternatif du transport déjà construit. Le noyau lit la retenue, le feuillet, l'orientation, la vacance et l'holonomie ; ses paramètres sont figés avant l'examen des angles externes.

\paragraph{Domaine et codomaine des réalisations CKM et PMNS}
La matrice rationnelle CKM, les repères \(F_\ell,F_\nu\), le transport \(O\), la phase \(\Phi_L\), l'unitarité de \(U_{\rm HMT}\) et le transport complet \(U_{\rm PMNS}=F_\ell^*F_\nu\) sont des résultats internes exacts sous les primitives déclarées. CKM constitue un transport physique une fois fixées les bases de saveur et de masse. La comparaison réfute uniquement la phase PMNS de rang un ; elle ne diminue pas la portée de la construction spectrale complète. La naturalité et le critère \(\det G\ne0\) appartiennent au certificat alternatif du registre et restent distincts de l'existence et de l'unitarité du transport PMNS.
